\begin{abstract}
\noindent
This document explores the profound cosmological and relativistic implications of the Navier--Stokes geometric phase topology. Building upon the rigorous formalization of the Bivariant Lie Algebra and the Yang-Mills Topological Equivalence, we demonstrate that the geometric saturation of a fluid's phase induces a strict dimensional collapse (3D $\to$ 2D $\to$ 1D). This topological strangulation generates a pure bivariant tension $\tau$ that acts as a fundamental field stress. We present the Quantum Astrophysical Fluid Topology (QAFT) axioms, demonstrating how the coherence barrier $\chi$ modulates the spacetime metric $g^{\mu\nu}_{\text{eff}}$, yielding an Extended General Relativity framework where fluid dynamics dictate metric geometry without singularities.
\end{abstract}

\maketitle

% ============================================================
% PREFACE
% ============================================================
\section*{Preface}
\addcontentsline{toc}{section}{Preface}

The mathematical resolution of the global regularity problem for the three-dimensional incompressible Navier--Stokes equations revealed an unexpected equivalence: the macroscopic fluid dynamics of the homochiral braking are mathematically isomorphic to the non-abelian gauge theory of Yang-Mills. This equivalence implies that the mechanisms preventing finite-time blow-up in classical fluids are topologically identical to the mass gap generating mechanisms in quantum vacuums.

This monograph, the second in our formal series, elevates these findings from fluid mechanics to General Relativity. By treating the triadic tensor as a strict Lie bracket and leveraging the Lean 4 formalization of the Saturn Symplectic Invariant, we construct the foundations of Functional Gravity. We show that phase saturation forces energy to bleed transversally into stable Yang-Mills rings, which in turn acts as an effective optical lens for the metric tensor.

The document is organized as follows:
Part I recalls the essential algebraic identities of the helicoidal triadic tensor and the Bivariant Lie algebra.
Part II introduces the dimensional phase transition, the pure bivariant tension $\tau$, and the Extended General Relativity equation.
